\subsection*{Intended Learning Outcomes}

The student should be able to:
\begin{enumerate}
  \item Chapter 3
    \begin{itemize}
      \item know and understand the definition of $i\sigma$-separation;
      \item check if sets are $i\sigma$-separated in given graphs;
    \end{itemize}
  \item Chapter 6
    \begin{itemize}
    \item know and understand the definition of structural causal model (SCM);
    \item know the differences between endogenous variables, exogenous input variables and exogenous random variables;
    \item understand the implicit assumptions made when using a SCM;
    \item know and understand the definitions of (potential) outcomes, solutions, their Markov kernels, and solution functions of SCMs;
    \item know how to sample from a SCM;
    \item know and understand the definitions of the different variants of hard interventions on SCMs;
    \item understand how soft interventions can be modeled in SCMs;
    \item understand how intervention variables can be used to model different interventions in the same causal model;
    \item understand how to compose and decompose SCMs;
    \item propose an appropriate SCM for a given real-world system, and to appropriately model interventions on the system using the model;
    \item know and understand the definition and consequences of unique solvability for SCMs;
    \item know and understand the definition and consequences of an SCM being simple;
    \item be able to check whether a given SCM is uniquely solvable/simple;
    \item construct observational and interventional Markov kernels for simple SCMs;
    \item understand the different equivalence relations for SCMs (see also Chapter 8);
    \item know and understand the definition of marginalization of SCMs;
    \item construct a marginalized SCM;
    \item know and apply the basic commutation and compatibility properties of basic operations on SCMs (intervening, marginalization, preservation of unique solvability / simplicity / equivalence);
    \item know and understand the definition of the graph of an SCM;
    \item be able to draw the graph of a given SCM, and to construct an SCM with a given graph;
    \item know and apply the basic compatibility relations between the operation that maps an SCM to a graph and other basic SCM operations;
    \item give a causal interpretation of the graph of an SCM (see Chapter 9);
    \item know the definition of acyclic SCM;
    \item understand and appreciate the differences between different causal model classes.
  \end{itemize}
  \item Chapter 7
    \begin{itemize}
      \item know and understand the definitions of acyclification of a graph and of a simple SCM, and their relation;
      \item know the relation between $i\sigma$-separation and $id$-separation on an acyclification;
      \item know and understand the global Markov property for simple SCMs;
      \item apply the global Markov property for simple SCMs;
      \item know and understand the 3 do-calculus rules for simple SCMs;
      \item apply the 3 do-calculus rules for simple SCMs;
      \item state and apply the back-door covariate adjustment criterion and formula for simple SCMs.
    \end{itemize}
  \item Chapter 8
    \begin{itemize}
      \item construct a twin SCM;
      \item construct a twin graph;
      \item know and understand the basic commutation and compatibility properties of twinning with other basic operations on and properties of SCMs;
      \item be able to reason about counterfactuals by using the twin SCM;
      \item be aware of the empirical limitations of working with counterfactuals.
    \end{itemize}
  \item Chapter 9
    \begin{itemize}
      \item know and understand the definitions of indirect causal relations, direct causal relations, and common causes according to simple SCMs;
      \item understand the relation between having a common cause and being confounded;
      \item express the notion of ``no confounding'' in terms of potential outcomes;
      \item know how to detect the presence of such relations from Markov kernels induced by simple SCMs;
      \item deal formally with a situation in which only part of the variables in an SCM are observed, whereas the others are latent;
      \item formally model a randomized controlled trial with SCMs in various ways, and prove in the SCM formalism the main properties of RCTs;
      \item know and understand the definition of faithfulness;
      \item understand different ways of how faithfulness violations can occur;
      \item understand the LCD algorithm;
      \item prove the correctness of the LCD algorithm;
      \item understand the Y-structures approach to causal discovery;
      \item prove the correctness of the Y-structures approach to causal discovery;
      \item know how minimal separating/connecting sets relate to ancestral relations.
    \end{itemize}
  \item Chapter 10
    \begin{itemize}
      \item understand how to use the $G$ test for conditional independence;
      \item understand under which conditions the $G$ test is valid (at level $\alpha$);
      \item understand under which conditions the $G$ test is asymptotically consistent;
      \item implement the $G$ test in a computer program;
    \end{itemize}
  \item Chapter 11
    \begin{itemize}
      \item know and understand the definition of $\sigma$-inducing path;
      \item know and understand the definition of a DPAG;
      \item know which DMGs a DPAG represents;
      \item construct a DPAG representing a given DMG.
    \end{itemize}
\end{enumerate}
